\begin{ExerciseBlock}
\ExerciseHeading{04}{Ultraschall in Werkstoffen}{wave_propagation_ultrasound_in_solids}
Verwenden Sie für alle Teilaufgaben das Modell longitudinaler Dehnwellen in einem schlanken Stab.
\begin{enumerate}
\item In Blei betragen die Phasengeschwindigkeit \qty{1226}{\metre\per\second} und die Dichte \qty{11.3}{\gram\per\centi\metre\cubed}. Wie groß ist der Elastizitätsmodul?
\item Bestimmen Sie die Phasengeschwindigkeit in Stahl mit $E=\qty{2.1e5}{\newton\per\milli\metre\squared}$ und $\rho=\qty{7.7}{\gram\per\centi\metre\cubed}$.
\item Ein zylindrischer Kupferstab ($\rho=\qty{8.92}{\gram\per\centi\metre\cubed}$) wird in Längsrichtung mit der Gewichtskraft einer Masse von \qty{40}{\kilogram} belastet. Ein Schallimpuls benötigt \qty{0.80}{\milli\second}, um den \qty{3.0}{\metre} langen Stab mit Durchmesser \qty{5}{\milli\metre} axial zu durchlaufen. Berechnen Sie den Betrag der Längenänderung $\Delta l$.
\end{enumerate}
\end{ExerciseBlock}
\UseExerciseSolution{wave_propagation_ultrasound_in_solids}
